% Generated draft: requires mathematical and source review.
\documentclass[11pt]{article}
\usepackage[margin=25mm]{geometry}
\usepackage{fontspec}
\setmainfont{Latin Modern Roman}
\usepackage{amsmath,amssymb,mathtools,unicode-math}
\setmathfont{Latin Modern Math}
\usepackage{enumitem,needspace,xurl,hyperref}
\setlength{\parindent}{0pt}
\setlength{\parskip}{5pt}
\setlength{\emergencystretch}{3em}
\newcommand{\N}{\mathbb{N}}
\newcommand{\Z}{\mathbb{Z}}
\newcommand{\Q}{\mathbb{Q}}
\newcommand{\R}{\mathbb{R}}
\newcommand{\C}{\mathbb{C}}
\newcommand{\F}{\mathbb{F}}
\newcommand{\A}{\mathbb{A}}
\newcommand{\PP}{\mathbb{P}}
\newcommand{\E}{\mathbb{E}}
\newcommand{\eps}{\varepsilon}
\newcommand{\dd}{\,\mathrm{d}}
\begin{document}
\section*{UnsolvedMath ID: 30002837}
\textbf{Source problem code:} catalogue entry
\subsection{Definitions and mathematical statement}
Let \(\mathrm{TMF}\) denote the spectrum of topological modular forms intended in the source, and write
\[
\mathrm{TMF}^{-n}(\mathrm{pt})=\pi_n(\mathrm{TMF}).
\]
The source defines \(\operatorname{Fer}(n)\) as \(n\) real chiral free fermions. This literal description requires \(n\) to be a nonnegative integer; the source specifies neither a further range nor an extension to other integers.

For a fixed such \(n\), let \(\mathcal M_n\) denote the intended moduli space of \(\operatorname{Fer}(n)\)-twisted extended two-dimensional field theories. The problem is conditional on the Stolz--Teichner identification cited in the source. Its precise meaning, including the geometric structures on the field theories and the model of the moduli space, is not supplied. In particular, the notation \(\mathrm{TMF}^{-n}\) in that identification is not specified as a representing space or as a cohomology group.

At the level of classes, suppose the intended identification induces
\[
\Phi_n:\pi_0(\mathcal M_n)\xrightarrow{\cong}
\mathrm{TMF}^{-n}(\mathrm{pt})=\pi_n(\mathrm{TMF}).
\]
Let \(\mathcal C_n\) denote the intended moduli space of conformal field theories admitting the required twisted extended field-theory interpretation, and suppose that interpretation gives a map \(j_n:\mathcal C_n\to\mathcal M_n\). These symbols describe the conditional formulation; the source does not construct either map or specify the admissible conformal field theories.

\textbf{Problem, conditional component-level formulation.} Determine the subset
\[
I_n=\operatorname{im}\!\left(
\pi_0(\mathcal C_n)\xrightarrow{(j_n)_*}
\pi_0(\mathcal M_n)\xrightarrow{\Phi_n}
\pi_n(\mathrm{TMF})
\right)
\]
represented by conformal field theories, and determine which classes regarded as interesting in the source can be obtained in this way. The source does not define ``interesting'' or name target classes, so it does not specify a precise existence assertion beyond this image question. It also asks how much of the field-theory side conformal field theories detect; whether that asks only about components or also about higher homotopy or other moduli-space information remains unresolved. No surjectivity, nonvanishing, or subgroup property of \(I_n\) is asserted.
\subsection{Short English statement}
Assuming the intended Stolz--Teichner correspondence, determine which classes in \(\pi_n(\mathrm{TMF})\) arise from conformal field theories with the required \(\operatorname{Fer}(n)\) twist and extension. The source leaves the admissible theories, the scope of detection, and the meaning of ``interesting classes'' unspecified. This is a conditional statement-only draft requiring human mathematical review.
\subsection{Sources}
\begin{itemize}
\item Local UnsolvedMath registry, permanent ID \(30002837\), title: \emph{Conformal Field Theories Representing TMF Classes}; exact source problem code: \texttt{catalogue entry}. The supplied record attributes the problem to Oberwolfach Reports workshop discussions and names no individual proposer.
\item Supplied report citation: \emph{Subfactors and Conformal Field Theory} (2016), Oberwolfach Reports, DOI \texttt{10.4171/owr/2015/16}, pp.\ 849--926. \url{https://doi.org/10.4171/owr/2015/16}. The registry separately records proposed year 2015; the relationship between these year labels has not been independently checked.
\item Supplied literature source, labeled ``Douglas--Henriques, conformal nets and TMF'': \url{https://arxiv.org/abs/1103.4187}.
\item Supplied literature source, labeled ``Henriques--Douglas, Stolz--Teichner cocycles'': \url{https://doi.org/10.4171/owr/2015/16}.
\item Literature assessment dated exactly \texttt{2026-08-21}: the supplied curation labels the problem open and reports that conformal-net and Stolz--Teichner constructions provide candidates and low-degree evidence, but that no general theorem realizes the requested nontrivial classes in \(\pi_n(\mathrm{TMF})\) by conformal field theories. This is supplied dated curation, not an independently verified assessment. The background explicitly requests independent verification of time-sensitive status and cited claims.
\end{itemize}
\end{document}
